% ####################################################################
% Basic configuration of Beamera class and Jagiellonian theme
% ####################################################################
\RequirePackage[l2tabu, orthodox]{nag}



\ifx\PresentationStyle\notset
  \def\PresentationStyle{dark}
\fi



% Options: t - align frame text to the top.
\documentclass[10pt,t]{beamer}
\mode<presentation>
\usetheme[style=\PresentationStyle]{jagiellonian}




% ------------------------------------------------------------------------------------
% Procesing configuration files of Jagiellonian theme located
% in the directory "preambule"
% ------------------------------------------------------------------------------------
% Configuration for polish language
% Need description
\usepackage[polish]{babel}
% Need description
\usepackage[MeX]{polski}



% ------------------------------
% Better support of polish chars in technical parts of PDF
% ------------------------------
\hypersetup{pdfencoding=auto,psdextra}

% Package "textpos" give as enviroment "textblock" which is very usefull in
% arranging text on slides.

% This is standard configuration of "textpos"
\usepackage[overlay,absolute]{textpos}

% If you need to see bounds of "textblock's" comment line above and uncomment
% one below.

% Caution! When showboxes option is on significant ammunt of space is add
% to the top of textblock and as such, everyting put in them gone down.
% We need to check how to remove this bug.

% \usepackage[showboxes,overlay,absolute]{textpos}



% Setting scale length for package "textpos"
\setlength{\TPHorizModule}{10mm}
\setlength{\TPVertModule}{\TPHorizModule}


% ---------------------------------------
% Packages written for lectures "Geometria 3D dla twórców gier wideo"
% ---------------------------------------
% \usepackage{./Geometry3DPackages/Geometry3D}
% \usepackage{./Geometry3DPackages/Geometry3DIndices}
% \usepackage{./Geometry3DPackages/Geometry3DTikZStyle}
% \usepackage{./ProgramowanieSymulacjiFizykiPaczki/ProgramowanieSymulacjiFizykiTikZStyle}
% \usepackage{./Geometry3DPackages/mathcommands}


% ---------------------------------------
% Importing TikZ's libraries
% ---------------------------------------
\usetikzlibrary{arrows.meta}
\usetikzlibrary{positioning}





% % Configuration package "bm" that need for making bold symbols
\newcommand{\bmmax}{0}
\newcommand{\hmmax}{0}
\usepackage{bm}




% ---------------------------------------
% Packages for scientific texts
% ---------------------------------------
% \let\lll\undefined  % Sometimes you must use this line to allow
% "amsmath" package to works with packages with packages for polish
% languge imported
% /preambul/LanguageSettings/JagiellonianPolishLanguageSettings.tex.
% This comments (probably) removes polish letter Ł.
\usepackage{amsmath}  % Packages from American Mathematical Society (AMS)
\usepackage{amssymb}
\usepackage{amscd}
\usepackage{amsthm}
\usepackage{siunitx}  % Package for typsetting SI units.
\usepackage{upgreek}  % Better looking greek letters.
% Example of using upgreek: pi = \uppi


\usepackage{calrsfs}  % Zmienia czcionkę kaligraficzną w \mathcal
% na ładniejszą. Może w innych miejscach robi to samo, ale o tym nic
% nie wiem.







% !!!!!!!!!!!!!!!!!!!!!!!!!!!!!!
% !!!!!!!!!!!!!!!!!!!!!!!!!!!!!!
% EVIL STUFF
\if\JUlogotitle1
\edef\LogoJUPath{LogoJU_\JUlogoLang/LogoJU_\JUlogoShape_\JUlogoColor.pdf}
\titlegraphic{\hfill\includegraphics[scale=0.22]
{./JagiellonianPictures/\LogoJUPath}}
\fi

% ---------------------------------------
% Commands for handling colors
% ---------------------------------------


% Command for setting normal text color for some text in math modestyle
% Text color depend on used style of Jagiellonian

% Beamer version of command
\newcommand{\TextWithNormalTextColor}[1]{%
  {\color{jNormalTextFGColor}
    \setbeamercolor{math text}{fg=jNormalTextFGColor} {#1}}
}

% Article and similar classes version of command
% \newcommand{\TextWithNormalTextColor}[1]{%
%   {\color{jNormalTextsFGColor} {#1}}
% }



% Beamer version of command
\newcommand{\NormalTextInMathMode}[1]{%
  {\color{jNormalTextFGColor}
    \setbeamercolor{math text}{fg=jNormalTextFGColor} \text{#1}}
}


% Article and similar classes version of command
% \newcommand{\NormalTextInMathMode}[1]{%
%   {\color{jNormalTextsFGColor} \text{#1}}
% }




% Command that sets color of some mathematical text to the same color
% that has normal text in header (?)

% Beamer version of the command
\newcommand{\MathTextFrametitleFGColor}[1]{%
  {\color{jFrametitleFGColor}
    \setbeamercolor{math text}{fg=jFrametitleFGColor} #1}
}

% Article and similar classes version of the command
% \newcommand{\MathTextWhiteColor}[1]{{\color{jFrametitleFGColor} #1}}





% Command for setting color of alert text for some text in math modestyle

% Beamer version of the command
\newcommand{\MathTextAlertColor}[1]{%
  {\color{jOrange} \setbeamercolor{math text}{fg=jOrange} #1}
}

% Article and similar classes version of the command
% \newcommand{\MathTextAlertColor}[1]{{\color{jOrange} #1}}





% Command that allow you to sets chosen color as the color of some text into
% math mode. Due to some nuances in the way that Beamer handle colors
% it not work in all cases. We hope that in the future we will improve it.

% Beamer version of the command
\newcommand{\SetMathTextColor}[2]{%
  {\color{#1} \setbeamercolor{math text}{fg=#1} #2}
}


% Article and similar classes version of the command
% \newcommand{\SetMathTextColor}[2]{{\color{#1} #2}}










% ---------------------------------------
% Commands for few special slides
% ---------------------------------------
\newcommand{\EndingSlide}[1]{%
  \begin{frame}[standout]

    \begingroup

    \color{jFrametitleFGColor}

    #1

    \endgroup

  \end{frame}
}










% ---------------------------------------
% Commands for setting background pictures for some slides
% ---------------------------------------
\newcommand{\TitleBackgroundPicture}
{./JagiellonianPictures/Backgrounds/LajkonikDark.png}
\newcommand{\SectionBackgroundPicture}
{./JagiellonianPictures/Backgrounds/LajkonikLight.png}



\newcommand{\TitleSlideWithPicture}{%
  \begingroup

  \usebackgroundtemplate{%
    \includegraphics[height=\paperheight]{\TitleBackgroundPicture}}

  \maketitle

  \endgroup
}





\newcommand{\SectionSlideWithPicture}[1]{%
  \begingroup

  \usebackgroundtemplate{%
    \includegraphics[height=\paperheight]{\SectionBackgroundPicture}}

  \setbeamercolor{titlelike}{fg=normal text.fg}

  \section{#1}

  \endgroup
}










% ---------------------------------------
% Commands for lectures "Geometria 3D dla twórców gier wideo"
% Polish version
% ---------------------------------------
% Komendy teraz wykomentowane były potrzebne, gdy loga były na niebieskim
% tle, nie na białym. A są na białym bo tego chcieli w biurze projektu.
% \newcommand{\FundingLogoWhitePicturePL}
% {./PresentationPictures/CommonPictures/logotypFundusze_biale_bez_tla2.pdf}
\newcommand{\FundingLogoColorPicturePL}
{./PresentationPictures/CommonPictures/European_Funds_color_PL.pdf}
% \newcommand{\EULogoWhitePicturePL}
% {./PresentationPictures/CommonPictures/logotypUE_biale_bez_tla2.pdf}
\newcommand{\EUSocialFundLogoColorPicturePL}
{./PresentationPictures/CommonPictures/EU_Social_Fund_color_PL.pdf}
% \newcommand{\ZintegrUJLogoWhitePicturePL}
% {./PresentationPictures/CommonPictures/zintegruj-logo-white.pdf}
\newcommand{\ZintegrUJLogoColorPicturePL}
{./PresentationPictures/CommonPictures/ZintegrUJ_color.pdf}
\newcommand{\JULogoColorPicturePL}
{./JagiellonianPictures/LogoJU_PL/LogoJU_A_color.pdf}





\newcommand{\GeometryThreeDSpecialBeginningSlidePL}{%
  \begin{frame}[standout]

    \begin{textblock}{11}(1,0.7)

      \begin{flushleft}

        \mdseries

        \footnotesize

        \color{jFrametitleFGColor}

        Materiał powstał w ramach projektu współfinansowanego ze środków
        Unii Europejskiej w ramach Europejskiego Funduszu Społecznego
        POWR.03.05.00-00-Z309/17-00.

      \end{flushleft}

    \end{textblock}





    \begin{textblock}{10}(0,2.2)

      \tikz \fill[color=jBackgroundStyleLight] (0,0) rectangle (12.8,-1.5);

    \end{textblock}


    \begin{textblock}{3.2}(1,2.45)

      \includegraphics[scale=0.3]{\FundingLogoColorPicturePL}

    \end{textblock}


    \begin{textblock}{2.5}(3.7,2.5)

      \includegraphics[scale=0.2]{\JULogoColorPicturePL}

    \end{textblock}


    \begin{textblock}{2.5}(6,2.4)

      \includegraphics[scale=0.1]{\ZintegrUJLogoColorPicturePL}

    \end{textblock}


    \begin{textblock}{4.2}(8.4,2.6)

      \includegraphics[scale=0.3]{\EUSocialFundLogoColorPicturePL}

    \end{textblock}

  \end{frame}
}



\newcommand{\GeometryThreeDTwoSpecialBeginningSlidesPL}{%
  \begin{frame}[standout]

    \begin{textblock}{11}(1,0.7)

      \begin{flushleft}

        \mdseries

        \footnotesize

        \color{jFrametitleFGColor}

        Materiał powstał w ramach projektu współfinansowanego ze środków
        Unii Europejskiej w ramach Europejskiego Funduszu Społecznego
        POWR.03.05.00-00-Z309/17-00.

      \end{flushleft}

    \end{textblock}





    \begin{textblock}{10}(0,2.2)

      \tikz \fill[color=jBackgroundStyleLight] (0,0) rectangle (12.8,-1.5);

    \end{textblock}


    \begin{textblock}{3.2}(1,2.45)

      \includegraphics[scale=0.3]{\FundingLogoColorPicturePL}

    \end{textblock}


    \begin{textblock}{2.5}(3.7,2.5)

      \includegraphics[scale=0.2]{\JULogoColorPicturePL}

    \end{textblock}


    \begin{textblock}{2.5}(6,2.4)

      \includegraphics[scale=0.1]{\ZintegrUJLogoColorPicturePL}

    \end{textblock}


    \begin{textblock}{4.2}(8.4,2.6)

      \includegraphics[scale=0.3]{\EUSocialFundLogoColorPicturePL}

    \end{textblock}

  \end{frame}





  \TitleSlideWithPicture
}



\newcommand{\GeometryThreeDSpecialEndingSlidePL}{%
  \begin{frame}[standout]

    \begin{textblock}{11}(1,0.7)

      \begin{flushleft}

        \mdseries

        \footnotesize

        \color{jFrametitleFGColor}

        Materiał powstał w ramach projektu współfinansowanego ze środków
        Unii Europejskiej w~ramach Europejskiego Funduszu Społecznego
        POWR.03.05.00-00-Z309/17-00.

      \end{flushleft}

    \end{textblock}





    \begin{textblock}{10}(0,2.2)

      \tikz \fill[color=jBackgroundStyleLight] (0,0) rectangle (12.8,-1.5);

    \end{textblock}


    \begin{textblock}{3.2}(1,2.45)

      \includegraphics[scale=0.3]{\FundingLogoColorPicturePL}

    \end{textblock}


    \begin{textblock}{2.5}(3.7,2.5)

      \includegraphics[scale=0.2]{\JULogoColorPicturePL}

    \end{textblock}


    \begin{textblock}{2.5}(6,2.4)

      \includegraphics[scale=0.1]{\ZintegrUJLogoColorPicturePL}

    \end{textblock}


    \begin{textblock}{4.2}(8.4,2.6)

      \includegraphics[scale=0.3]{\EUSocialFundLogoColorPicturePL}

    \end{textblock}





    \begin{textblock}{11}(1,4)

      \begin{flushleft}

        \mdseries

        \footnotesize

        \RaggedRight

        \color{jFrametitleFGColor}

        Treść niniejszego wykładu jest udostępniona na~licencji
        Creative Commons (\textsc{cc}), z~uzna\-niem autorstwa
        (\textsc{by}) oraz udostępnianiem na tych samych warunkach
        (\textsc{sa}). Rysunki i~wy\-kresy zawarte w~wykładzie są
        autorstwa dr.~hab.~Pawła Węgrzyna et~al. i~są dostępne
        na tej samej licencji, o~ile nie wskazano inaczej.
        W~prezentacji wykorzystano temat Beamera Jagiellonian,
        oparty na~temacie Metropolis Matthiasa Vogelgesanga,
        dostępnym na licencji \LaTeX{} Project Public License~1.3c
        pod adresem: \colorhref{https://github.com/matze/mtheme}
        {https://github.com/matze/mtheme}.

        Projekt typograficzny: Iwona Grabska-Gradzińska \\
        Skład: Kamil Ziemian;
        Korekta: Wojciech Palacz \\
        Modele: Dariusz Frymus, Kamil Nowakowski \\
        Rysunki i~wykresy: Kamil Ziemian, Paweł Węgrzyn, Wojciech Palacz

      \end{flushleft}

    \end{textblock}

  \end{frame}
}



\newcommand{\GeometryThreeDTwoSpecialEndingSlidesPL}[1]{%
  \begin{frame}[standout]


    \begin{textblock}{11}(1,0.7)

      \begin{flushleft}

        \mdseries

        \footnotesize

        \color{jFrametitleFGColor}

        Materiał powstał w ramach projektu współfinansowanego ze środków
        Unii Europejskiej w~ramach Europejskiego Funduszu Społecznego
        POWR.03.05.00-00-Z309/17-00.

      \end{flushleft}

    \end{textblock}





    \begin{textblock}{10}(0,2.2)

      \tikz \fill[color=jBackgroundStyleLight] (0,0) rectangle (12.8,-1.5);

    \end{textblock}


    \begin{textblock}{3.2}(1,2.45)

      \includegraphics[scale=0.3]{\FundingLogoColorPicturePL}

    \end{textblock}


    \begin{textblock}{2.5}(3.7,2.5)

      \includegraphics[scale=0.2]{\JULogoColorPicturePL}

    \end{textblock}


    \begin{textblock}{2.5}(6,2.4)

      \includegraphics[scale=0.1]{\ZintegrUJLogoColorPicturePL}

    \end{textblock}


    \begin{textblock}{4.2}(8.4,2.6)

      \includegraphics[scale=0.3]{\EUSocialFundLogoColorPicturePL}

    \end{textblock}





    \begin{textblock}{11}(1,4)

      \begin{flushleft}

        \mdseries

        \footnotesize

        \RaggedRight

        \color{jFrametitleFGColor}

        Treść niniejszego wykładu jest udostępniona na~licencji
        Creative Commons (\textsc{cc}), z~uzna\-niem autorstwa
        (\textsc{by}) oraz udostępnianiem na tych samych warunkach
        (\textsc{sa}). Rysunki i~wy\-kresy zawarte w~wykładzie są
        autorstwa dr.~hab.~Pawła Węgrzyna et~al. i~są dostępne
        na tej samej licencji, o~ile nie wskazano inaczej.
        W~prezentacji wykorzystano temat Beamera Jagiellonian,
        oparty na~temacie Metropolis Matthiasa Vogelgesanga,
        dostępnym na licencji \LaTeX{} Project Public License~1.3c
        pod adresem: \colorhref{https://github.com/matze/mtheme}
        {https://github.com/matze/mtheme}.

        Projekt typograficzny: Iwona Grabska-Gradzińska \\
        Skład: Kamil Ziemian;
        Korekta: Wojciech Palacz \\
        Modele: Dariusz Frymus, Kamil Nowakowski \\
        Rysunki i~wykresy: Kamil Ziemian, Paweł Węgrzyn, Wojciech Palacz

      \end{flushleft}

    \end{textblock}

  \end{frame}





  \begin{frame}[standout]

    \begingroup

    \color{jFrametitleFGColor}

    #1

    \endgroup

  \end{frame}
}



\newcommand{\GeometryThreeDSpecialEndingSlideVideoPL}{%
  \begin{frame}[standout]

    \begin{textblock}{11}(1,0.7)

      \begin{flushleft}

        \mdseries

        \footnotesize

        \color{jFrametitleFGColor}

        Materiał powstał w ramach projektu współfinansowanego ze środków
        Unii Europejskiej w~ramach Europejskiego Funduszu Społecznego
        POWR.03.05.00-00-Z309/17-00.

      \end{flushleft}

    \end{textblock}





    \begin{textblock}{10}(0,2.2)

      \tikz \fill[color=jBackgroundStyleLight] (0,0) rectangle (12.8,-1.5);

    \end{textblock}


    \begin{textblock}{3.2}(1,2.45)

      \includegraphics[scale=0.3]{\FundingLogoColorPicturePL}

    \end{textblock}


    \begin{textblock}{2.5}(3.7,2.5)

      \includegraphics[scale=0.2]{\JULogoColorPicturePL}

    \end{textblock}


    \begin{textblock}{2.5}(6,2.4)

      \includegraphics[scale=0.1]{\ZintegrUJLogoColorPicturePL}

    \end{textblock}


    \begin{textblock}{4.2}(8.4,2.6)

      \includegraphics[scale=0.3]{\EUSocialFundLogoColorPicturePL}

    \end{textblock}





    \begin{textblock}{11}(1,4)

      \begin{flushleft}

        \mdseries

        \footnotesize

        \RaggedRight

        \color{jFrametitleFGColor}

        Treść niniejszego wykładu jest udostępniona na~licencji
        Creative Commons (\textsc{cc}), z~uzna\-niem autorstwa
        (\textsc{by}) oraz udostępnianiem na tych samych warunkach
        (\textsc{sa}). Rysunki i~wy\-kresy zawarte w~wykładzie są
        autorstwa dr.~hab.~Pawła Węgrzyna et~al. i~są dostępne
        na tej samej licencji, o~ile nie wskazano inaczej.
        W~prezentacji wykorzystano temat Beamera Jagiellonian,
        oparty na~temacie Metropolis Matthiasa Vogelgesanga,
        dostępnym na licencji \LaTeX{} Project Public License~1.3c
        pod adresem: \colorhref{https://github.com/matze/mtheme}
        {https://github.com/matze/mtheme}.

        Projekt typograficzny: Iwona Grabska-Gradzińska;
        Skład: Kamil Ziemian \\
        Korekta: Wojciech Palacz;
        Modele: Dariusz Frymus, Kamil Nowakowski \\
        Rysunki i~wykresy: Kamil Ziemian, Paweł Węgrzyn, Wojciech Palacz \\
        Montaż: Agencja Filmowa Film \& Television Production~-- Zbigniew
        Masklak

      \end{flushleft}

    \end{textblock}

  \end{frame}
}





\newcommand{\GeometryThreeDTwoSpecialEndingSlidesVideoPL}[1]{%
  \begin{frame}[standout]

    \begin{textblock}{11}(1,0.7)

      \begin{flushleft}

        \mdseries

        \footnotesize

        \color{jFrametitleFGColor}

        Materiał powstał w ramach projektu współfinansowanego ze środków
        Unii Europejskiej w~ramach Europejskiego Funduszu Społecznego
        POWR.03.05.00-00-Z309/17-00.

      \end{flushleft}

    \end{textblock}





    \begin{textblock}{10}(0,2.2)

      \tikz \fill[color=jBackgroundStyleLight] (0,0) rectangle (12.8,-1.5);

    \end{textblock}


    \begin{textblock}{3.2}(1,2.45)

      \includegraphics[scale=0.3]{\FundingLogoColorPicturePL}

    \end{textblock}


    \begin{textblock}{2.5}(3.7,2.5)

      \includegraphics[scale=0.2]{\JULogoColorPicturePL}

    \end{textblock}


    \begin{textblock}{2.5}(6,2.4)

      \includegraphics[scale=0.1]{\ZintegrUJLogoColorPicturePL}

    \end{textblock}


    \begin{textblock}{4.2}(8.4,2.6)

      \includegraphics[scale=0.3]{\EUSocialFundLogoColorPicturePL}

    \end{textblock}





    \begin{textblock}{11}(1,4)

      \begin{flushleft}

        \mdseries

        \footnotesize

        \RaggedRight

        \color{jFrametitleFGColor}

        Treść niniejszego wykładu jest udostępniona na~licencji
        Creative Commons (\textsc{cc}), z~uzna\-niem autorstwa
        (\textsc{by}) oraz udostępnianiem na tych samych warunkach
        (\textsc{sa}). Rysunki i~wy\-kresy zawarte w~wykładzie są
        autorstwa dr.~hab.~Pawła Węgrzyna et~al. i~są dostępne
        na tej samej licencji, o~ile nie wskazano inaczej.
        W~prezentacji wykorzystano temat Beamera Jagiellonian,
        oparty na~temacie Metropolis Matthiasa Vogelgesanga,
        dostępnym na licencji \LaTeX{} Project Public License~1.3c
        pod adresem: \colorhref{https://github.com/matze/mtheme}
        {https://github.com/matze/mtheme}.

        Projekt typograficzny: Iwona Grabska-Gradzińska;
        Skład: Kamil Ziemian \\
        Korekta: Wojciech Palacz;
        Modele: Dariusz Frymus, Kamil Nowakowski \\
        Rysunki i~wykresy: Kamil Ziemian, Paweł Węgrzyn, Wojciech Palacz \\
        Montaż: Agencja Filmowa Film \& Television Production~-- Zbigniew
        Masklak

      \end{flushleft}

    \end{textblock}

  \end{frame}





  \begin{frame}[standout]


    \begingroup

    \color{jFrametitleFGColor}

    #1

    \endgroup

  \end{frame}
}










% ---------------------------------------
% Commands for lectures "Geometria 3D dla twórców gier wideo"
% English version
% ---------------------------------------
% \newcommand{\FundingLogoWhitePictureEN}
% {./PresentationPictures/CommonPictures/logotypFundusze_biale_bez_tla2.pdf}
\newcommand{\FundingLogoColorPictureEN}
{./PresentationPictures/CommonPictures/European_Funds_color_EN.pdf}
% \newcommand{\EULogoWhitePictureEN}
% {./PresentationPictures/CommonPictures/logotypUE_biale_bez_tla2.pdf}
\newcommand{\EUSocialFundLogoColorPictureEN}
{./PresentationPictures/CommonPictures/EU_Social_Fund_color_EN.pdf}
% \newcommand{\ZintegrUJLogoWhitePictureEN}
% {./PresentationPictures/CommonPictures/zintegruj-logo-white.pdf}
\newcommand{\ZintegrUJLogoColorPictureEN}
{./PresentationPictures/CommonPictures/ZintegrUJ_color.pdf}
\newcommand{\JULogoColorPictureEN}
{./JagiellonianPictures/LogoJU_EN/LogoJU_A_color.pdf}



\newcommand{\GeometryThreeDSpecialBeginningSlideEN}{%
  \begin{frame}[standout]

    \begin{textblock}{11}(1,0.7)

      \begin{flushleft}

        \mdseries

        \footnotesize

        \color{jFrametitleFGColor}

        This content was created as part of a project co-financed by the
        European Union within the framework of the European Social Fund
        POWR.03.05.00-00-Z309/17-00.

      \end{flushleft}

    \end{textblock}





    \begin{textblock}{10}(0,2.2)

      \tikz \fill[color=jBackgroundStyleLight] (0,0) rectangle (12.8,-1.5);

    \end{textblock}


    \begin{textblock}{3.2}(0.7,2.45)

      \includegraphics[scale=0.3]{\FundingLogoColorPictureEN}

    \end{textblock}


    \begin{textblock}{2.5}(4.15,2.5)

      \includegraphics[scale=0.2]{\JULogoColorPictureEN}

    \end{textblock}


    \begin{textblock}{2.5}(6.35,2.4)

      \includegraphics[scale=0.1]{\ZintegrUJLogoColorPictureEN}

    \end{textblock}


    \begin{textblock}{4.2}(8.4,2.6)

      \includegraphics[scale=0.3]{\EUSocialFundLogoColorPictureEN}

    \end{textblock}

  \end{frame}
}



\newcommand{\GeometryThreeDTwoSpecialBeginningSlidesEN}{%
  \begin{frame}[standout]

    \begin{textblock}{11}(1,0.7)

      \begin{flushleft}

        \mdseries

        \footnotesize

        \color{jFrametitleFGColor}

        This content was created as part of a project co-financed by the
        European Union within the framework of the European Social Fund
        POWR.03.05.00-00-Z309/17-00.

      \end{flushleft}

    \end{textblock}





    \begin{textblock}{10}(0,2.2)

      \tikz \fill[color=jBackgroundStyleLight] (0,0) rectangle (12.8,-1.5);

    \end{textblock}


    \begin{textblock}{3.2}(0.7,2.45)

      \includegraphics[scale=0.3]{\FundingLogoColorPictureEN}

    \end{textblock}


    \begin{textblock}{2.5}(4.15,2.5)

      \includegraphics[scale=0.2]{\JULogoColorPictureEN}

    \end{textblock}


    \begin{textblock}{2.5}(6.35,2.4)

      \includegraphics[scale=0.1]{\ZintegrUJLogoColorPictureEN}

    \end{textblock}


    \begin{textblock}{4.2}(8.4,2.6)

      \includegraphics[scale=0.3]{\EUSocialFundLogoColorPictureEN}

    \end{textblock}

  \end{frame}





  \TitleSlideWithPicture
}



\newcommand{\GeometryThreeDSpecialEndingSlideEN}{%
  \begin{frame}[standout]

    \begin{textblock}{11}(1,0.7)

      \begin{flushleft}

        \mdseries

        \footnotesize

        \color{jFrametitleFGColor}

        This content was created as part of a project co-financed by the
        European Union within the framework of the European Social Fund
        POWR.03.05.00-00-Z309/17-00.

      \end{flushleft}

    \end{textblock}





    \begin{textblock}{10}(0,2.2)

      \tikz \fill[color=jBackgroundStyleLight] (0,0) rectangle (12.8,-1.5);

    \end{textblock}


    \begin{textblock}{3.2}(0.7,2.45)

      \includegraphics[scale=0.3]{\FundingLogoColorPictureEN}

    \end{textblock}


    \begin{textblock}{2.5}(4.15,2.5)

      \includegraphics[scale=0.2]{\JULogoColorPictureEN}

    \end{textblock}


    \begin{textblock}{2.5}(6.35,2.4)

      \includegraphics[scale=0.1]{\ZintegrUJLogoColorPictureEN}

    \end{textblock}


    \begin{textblock}{4.2}(8.4,2.6)

      \includegraphics[scale=0.3]{\EUSocialFundLogoColorPictureEN}

    \end{textblock}





    \begin{textblock}{11}(1,4)

      \begin{flushleft}

        \mdseries

        \footnotesize

        \RaggedRight

        \color{jFrametitleFGColor}

        The content of this lecture is made available under a~Creative
        Commons licence (\textsc{cc}), giving the author the credits
        (\textsc{by}) and putting an obligation to share on the same terms
        (\textsc{sa}). Figures and diagrams included in the lecture are
        authored by Paweł Węgrzyn et~al., and are available under the same
        license unless indicated otherwise.\\ The presentation uses the
        Beamer Jagiellonian theme based on Matthias Vogelgesang’s
        Metropolis theme, available under license \LaTeX{} Project
        Public License~1.3c at: \colorhref{https://github.com/matze/mtheme}
        {https://github.com/matze/mtheme}.

        Typographic design: Iwona Grabska-Gradzińska \\
        \LaTeX{} Typesetting: Kamil Ziemian \\
        Proofreading: Wojciech Palacz,
        Monika Stawicka \\
        3D Models: Dariusz Frymus, Kamil Nowakowski \\
        Figures and charts: Kamil Ziemian, Paweł Węgrzyn, Wojciech Palacz

      \end{flushleft}

    \end{textblock}

  \end{frame}
}



\newcommand{\GeometryThreeDTwoSpecialEndingSlidesEN}[1]{%
  \begin{frame}[standout]


    \begin{textblock}{11}(1,0.7)

      \begin{flushleft}

        \mdseries

        \footnotesize

        \color{jFrametitleFGColor}

        This content was created as part of a project co-financed by the
        European Union within the framework of the European Social Fund
        POWR.03.05.00-00-Z309/17-00.

      \end{flushleft}

    \end{textblock}





    \begin{textblock}{10}(0,2.2)

      \tikz \fill[color=jBackgroundStyleLight] (0,0) rectangle (12.8,-1.5);

    \end{textblock}


    \begin{textblock}{3.2}(0.7,2.45)

      \includegraphics[scale=0.3]{\FundingLogoColorPictureEN}

    \end{textblock}


    \begin{textblock}{2.5}(4.15,2.5)

      \includegraphics[scale=0.2]{\JULogoColorPictureEN}

    \end{textblock}


    \begin{textblock}{2.5}(6.35,2.4)

      \includegraphics[scale=0.1]{\ZintegrUJLogoColorPictureEN}

    \end{textblock}


    \begin{textblock}{4.2}(8.4,2.6)

      \includegraphics[scale=0.3]{\EUSocialFundLogoColorPictureEN}

    \end{textblock}





    \begin{textblock}{11}(1,4)

      \begin{flushleft}

        \mdseries

        \footnotesize

        \RaggedRight

        \color{jFrametitleFGColor}

        The content of this lecture is made available under a~Creative
        Commons licence (\textsc{cc}), giving the author the credits
        (\textsc{by}) and putting an obligation to share on the same terms
        (\textsc{sa}). Figures and diagrams included in the lecture are
        authored by Paweł Węgrzyn et~al., and are available under the same
        license unless indicated otherwise.\\ The presentation uses the
        Beamer Jagiellonian theme based on Matthias Vogelgesang’s
        Metropolis theme, available under license \LaTeX{} Project
        Public License~1.3c at: \colorhref{https://github.com/matze/mtheme}
        {https://github.com/matze/mtheme}.

        Typographic design: Iwona Grabska-Gradzińska \\
        \LaTeX{} Typesetting: Kamil Ziemian \\
        Proofreading: Wojciech Palacz,
        Monika Stawicka \\
        3D Models: Dariusz Frymus, Kamil Nowakowski \\
        Figures and charts: Kamil Ziemian, Paweł Węgrzyn, Wojciech Palacz

      \end{flushleft}

    \end{textblock}

  \end{frame}





  \begin{frame}[standout]

    \begingroup

    \color{jFrametitleFGColor}

    #1

    \endgroup

  \end{frame}
}



\newcommand{\GeometryThreeDSpecialEndingSlideVideoVerOneEN}{%
  \begin{frame}[standout]

    \begin{textblock}{11}(1,0.7)

      \begin{flushleft}

        \mdseries

        \footnotesize

        \color{jFrametitleFGColor}

        This content was created as part of a project co-financed by the
        European Union within the framework of the European Social Fund
        POWR.03.05.00-00-Z309/17-00.

      \end{flushleft}

    \end{textblock}





    \begin{textblock}{10}(0,2.2)

      \tikz \fill[color=jBackgroundStyleLight] (0,0) rectangle (12.8,-1.5);

    \end{textblock}


    \begin{textblock}{3.2}(0.7,2.45)

      \includegraphics[scale=0.3]{\FundingLogoColorPictureEN}

    \end{textblock}


    \begin{textblock}{2.5}(4.15,2.5)

      \includegraphics[scale=0.2]{\JULogoColorPictureEN}

    \end{textblock}


    \begin{textblock}{2.5}(6.35,2.4)

      \includegraphics[scale=0.1]{\ZintegrUJLogoColorPictureEN}

    \end{textblock}


    \begin{textblock}{4.2}(8.4,2.6)

      \includegraphics[scale=0.3]{\EUSocialFundLogoColorPictureEN}

    \end{textblock}





    \begin{textblock}{11}(1,4)

      \begin{flushleft}

        \mdseries

        \footnotesize

        \RaggedRight

        \color{jFrametitleFGColor}

        The content of this lecture is made available under a Creative
        Commons licence (\textsc{cc}), giving the author the credits
        (\textsc{by}) and putting an obligation to share on the same terms
        (\textsc{sa}). Figures and diagrams included in the lecture are
        authored by Paweł Węgrzyn et~al., and are available under the same
        license unless indicated otherwise.\\ The presentation uses the
        Beamer Jagiellonian theme based on Matthias Vogelgesang’s
        Metropolis theme, available under license \LaTeX{} Project
        Public License~1.3c at: \colorhref{https://github.com/matze/mtheme}
        {https://github.com/matze/mtheme}.

        Typographic design: Iwona Grabska-Gradzińska;
        \LaTeX{} Typesetting: Kamil Ziemian \\
        Proofreading: Wojciech Palacz,
        Monika Stawicka \\
        3D Models: Dariusz Frymus, Kamil Nowakowski \\
        Figures and charts: Kamil Ziemian, Paweł Węgrzyn, Wojciech
        Palacz \\
        Film editing: Agencja Filmowa Film \& Television Production~--
        Zbigniew Masklak

      \end{flushleft}

    \end{textblock}

  \end{frame}
}



\newcommand{\GeometryThreeDSpecialEndingSlideVideoVerTwoEN}{%
  \begin{frame}[standout]

    \begin{textblock}{11}(1,0.7)

      \begin{flushleft}

        \mdseries

        \footnotesize

        \color{jFrametitleFGColor}

        This content was created as part of a project co-financed by the
        European Union within the framework of the European Social Fund
        POWR.03.05.00-00-Z309/17-00.

      \end{flushleft}

    \end{textblock}





    \begin{textblock}{10}(0,2.2)

      \tikz \fill[color=jBackgroundStyleLight] (0,0) rectangle (12.8,-1.5);

    \end{textblock}


    \begin{textblock}{3.2}(0.7,2.45)

      \includegraphics[scale=0.3]{\FundingLogoColorPictureEN}

    \end{textblock}


    \begin{textblock}{2.5}(4.15,2.5)

      \includegraphics[scale=0.2]{\JULogoColorPictureEN}

    \end{textblock}


    \begin{textblock}{2.5}(6.35,2.4)

      \includegraphics[scale=0.1]{\ZintegrUJLogoColorPictureEN}

    \end{textblock}


    \begin{textblock}{4.2}(8.4,2.6)

      \includegraphics[scale=0.3]{\EUSocialFundLogoColorPictureEN}

    \end{textblock}





    \begin{textblock}{11}(1,4)

      \begin{flushleft}

        \mdseries

        \footnotesize

        \RaggedRight

        \color{jFrametitleFGColor}

        The content of this lecture is made available under a Creative
        Commons licence (\textsc{cc}), giving the author the credits
        (\textsc{by}) and putting an obligation to share on the same terms
        (\textsc{sa}). Figures and diagrams included in the lecture are
        authored by Paweł Węgrzyn et~al., and are available under the same
        license unless indicated otherwise.\\ The presentation uses the
        Beamer Jagiellonian theme based on Matthias Vogelgesang’s
        Metropolis theme, available under license \LaTeX{} Project
        Public License~1.3c at: \colorhref{https://github.com/matze/mtheme}
        {https://github.com/matze/mtheme}.

        Typographic design: Iwona Grabska-Gradzińska;
        \LaTeX{} Typesetting: Kamil Ziemian \\
        Proofreading: Wojciech Palacz,
        Monika Stawicka \\
        3D Models: Dariusz Frymus, Kamil Nowakowski \\
        Figures and charts: Kamil Ziemian, Paweł Węgrzyn, Wojciech
        Palacz \\
        Film editing: IMAVI -- Joanna Kozakiewicz, Krzysztof Magda, Nikodem
        Frodyma

      \end{flushleft}

    \end{textblock}

  \end{frame}
}



\newcommand{\GeometryThreeDSpecialEndingSlideVideoVerThreeEN}{%
  \begin{frame}[standout]

    \begin{textblock}{11}(1,0.7)

      \begin{flushleft}

        \mdseries

        \footnotesize

        \color{jFrametitleFGColor}

        This content was created as part of a project co-financed by the
        European Union within the framework of the European Social Fund
        POWR.03.05.00-00-Z309/17-00.

      \end{flushleft}

    \end{textblock}





    \begin{textblock}{10}(0,2.2)

      \tikz \fill[color=jBackgroundStyleLight] (0,0) rectangle (12.8,-1.5);

    \end{textblock}


    \begin{textblock}{3.2}(0.7,2.45)

      \includegraphics[scale=0.3]{\FundingLogoColorPictureEN}

    \end{textblock}


    \begin{textblock}{2.5}(4.15,2.5)

      \includegraphics[scale=0.2]{\JULogoColorPictureEN}

    \end{textblock}


    \begin{textblock}{2.5}(6.35,2.4)

      \includegraphics[scale=0.1]{\ZintegrUJLogoColorPictureEN}

    \end{textblock}


    \begin{textblock}{4.2}(8.4,2.6)

      \includegraphics[scale=0.3]{\EUSocialFundLogoColorPictureEN}

    \end{textblock}





    \begin{textblock}{11}(1,4)

      \begin{flushleft}

        \mdseries

        \footnotesize

        \RaggedRight

        \color{jFrametitleFGColor}

        The content of this lecture is made available under a Creative
        Commons licence (\textsc{cc}), giving the author the credits
        (\textsc{by}) and putting an obligation to share on the same terms
        (\textsc{sa}). Figures and diagrams included in the lecture are
        authored by Paweł Węgrzyn et~al., and are available under the same
        license unless indicated otherwise.\\ The presentation uses the
        Beamer Jagiellonian theme based on Matthias Vogelgesang’s
        Metropolis theme, available under license \LaTeX{} Project
        Public License~1.3c at: \colorhref{https://github.com/matze/mtheme}
        {https://github.com/matze/mtheme}.

        Typographic design: Iwona Grabska-Gradzińska;
        \LaTeX{} Typesetting: Kamil Ziemian \\
        Proofreading: Wojciech Palacz,
        Monika Stawicka \\
        3D Models: Dariusz Frymus, Kamil Nowakowski \\
        Figures and charts: Kamil Ziemian, Paweł Węgrzyn, Wojciech
        Palacz \\
        Film editing: Agencja Filmowa Film \& Television Production~--
        Zbigniew Masklak \\
        Film editing: IMAVI -- Joanna Kozakiewicz, Krzysztof Magda, Nikodem
        Frodyma

      \end{flushleft}

    \end{textblock}

  \end{frame}
}



\newcommand{\GeometryThreeDTwoSpecialEndingSlidesVideoVerOneEN}[1]{%
  \begin{frame}[standout]

    \begin{textblock}{11}(1,0.7)

      \begin{flushleft}

        \mdseries

        \footnotesize

        \color{jFrametitleFGColor}

        This content was created as part of a project co-financed by the
        European Union within the framework of the European Social Fund
        POWR.03.05.00-00-Z309/17-00.

      \end{flushleft}

    \end{textblock}





    \begin{textblock}{10}(0,2.2)

      \tikz \fill[color=jBackgroundStyleLight] (0,0) rectangle (12.8,-1.5);

    \end{textblock}


    \begin{textblock}{3.2}(0.7,2.45)

      \includegraphics[scale=0.3]{\FundingLogoColorPictureEN}

    \end{textblock}


    \begin{textblock}{2.5}(4.15,2.5)

      \includegraphics[scale=0.2]{\JULogoColorPictureEN}

    \end{textblock}


    \begin{textblock}{2.5}(6.35,2.4)

      \includegraphics[scale=0.1]{\ZintegrUJLogoColorPictureEN}

    \end{textblock}


    \begin{textblock}{4.2}(8.4,2.6)

      \includegraphics[scale=0.3]{\EUSocialFundLogoColorPictureEN}

    \end{textblock}





    \begin{textblock}{11}(1,4)

      \begin{flushleft}

        \mdseries

        \footnotesize

        \RaggedRight

        \color{jFrametitleFGColor}

        The content of this lecture is made available under a Creative
        Commons licence (\textsc{cc}), giving the author the credits
        (\textsc{by}) and putting an obligation to share on the same terms
        (\textsc{sa}). Figures and diagrams included in the lecture are
        authored by Paweł Węgrzyn et~al., and are available under the same
        license unless indicated otherwise.\\ The presentation uses the
        Beamer Jagiellonian theme based on Matthias Vogelgesang’s
        Metropolis theme, available under license \LaTeX{} Project
        Public License~1.3c at: \colorhref{https://github.com/matze/mtheme}
        {https://github.com/matze/mtheme}.

        Typographic design: Iwona Grabska-Gradzińska;
        \LaTeX{} Typesetting: Kamil Ziemian \\
        Proofreading: Wojciech Palacz,
        Monika Stawicka \\
        3D Models: Dariusz Frymus, Kamil Nowakowski \\
        Figures and charts: Kamil Ziemian, Paweł Węgrzyn,
        Wojciech Palacz \\
        Film editing: Agencja Filmowa Film \& Television Production~--
        Zbigniew Masklak

      \end{flushleft}

    \end{textblock}

  \end{frame}





  \begin{frame}[standout]


    \begingroup

    \color{jFrametitleFGColor}

    #1

    \endgroup

  \end{frame}
}



\newcommand{\GeometryThreeDTwoSpecialEndingSlidesVideoVerTwoEN}[1]{%
  \begin{frame}[standout]

    \begin{textblock}{11}(1,0.7)

      \begin{flushleft}

        \mdseries

        \footnotesize

        \color{jFrametitleFGColor}

        This content was created as part of a project co-financed by the
        European Union within the framework of the European Social Fund
        POWR.03.05.00-00-Z309/17-00.

      \end{flushleft}

    \end{textblock}





    \begin{textblock}{10}(0,2.2)

      \tikz \fill[color=jBackgroundStyleLight] (0,0) rectangle (12.8,-1.5);

    \end{textblock}


    \begin{textblock}{3.2}(0.7,2.45)

      \includegraphics[scale=0.3]{\FundingLogoColorPictureEN}

    \end{textblock}


    \begin{textblock}{2.5}(4.15,2.5)

      \includegraphics[scale=0.2]{\JULogoColorPictureEN}

    \end{textblock}


    \begin{textblock}{2.5}(6.35,2.4)

      \includegraphics[scale=0.1]{\ZintegrUJLogoColorPictureEN}

    \end{textblock}


    \begin{textblock}{4.2}(8.4,2.6)

      \includegraphics[scale=0.3]{\EUSocialFundLogoColorPictureEN}

    \end{textblock}





    \begin{textblock}{11}(1,4)

      \begin{flushleft}

        \mdseries

        \footnotesize

        \RaggedRight

        \color{jFrametitleFGColor}

        The content of this lecture is made available under a Creative
        Commons licence (\textsc{cc}), giving the author the credits
        (\textsc{by}) and putting an obligation to share on the same terms
        (\textsc{sa}). Figures and diagrams included in the lecture are
        authored by Paweł Węgrzyn et~al., and are available under the same
        license unless indicated otherwise.\\ The presentation uses the
        Beamer Jagiellonian theme based on Matthias Vogelgesang’s
        Metropolis theme, available under license \LaTeX{} Project
        Public License~1.3c at: \colorhref{https://github.com/matze/mtheme}
        {https://github.com/matze/mtheme}.

        Typographic design: Iwona Grabska-Gradzińska;
        \LaTeX{} Typesetting: Kamil Ziemian \\
        Proofreading: Wojciech Palacz,
        Monika Stawicka \\
        3D Models: Dariusz Frymus, Kamil Nowakowski \\
        Figures and charts: Kamil Ziemian, Paweł Węgrzyn,
        Wojciech Palacz \\
        Film editing: IMAVI -- Joanna Kozakiewicz, Krzysztof Magda, Nikodem
        Frodyma

      \end{flushleft}

    \end{textblock}

  \end{frame}





  \begin{frame}[standout]


    \begingroup

    \color{jFrametitleFGColor}

    #1

    \endgroup

  \end{frame}
}



\newcommand{\GeometryThreeDTwoSpecialEndingSlidesVideoVerThreeEN}[1]{%
  \begin{frame}[standout]

    \begin{textblock}{11}(1,0.7)

      \begin{flushleft}

        \mdseries

        \footnotesize

        \color{jFrametitleFGColor}

        This content was created as part of a project co-financed by the
        European Union within the framework of the European Social Fund
        POWR.03.05.00-00-Z309/17-00.

      \end{flushleft}

    \end{textblock}





    \begin{textblock}{10}(0,2.2)

      \tikz \fill[color=jBackgroundStyleLight] (0,0) rectangle (12.8,-1.5);

    \end{textblock}


    \begin{textblock}{3.2}(0.7,2.45)

      \includegraphics[scale=0.3]{\FundingLogoColorPictureEN}

    \end{textblock}


    \begin{textblock}{2.5}(4.15,2.5)

      \includegraphics[scale=0.2]{\JULogoColorPictureEN}

    \end{textblock}


    \begin{textblock}{2.5}(6.35,2.4)

      \includegraphics[scale=0.1]{\ZintegrUJLogoColorPictureEN}

    \end{textblock}


    \begin{textblock}{4.2}(8.4,2.6)

      \includegraphics[scale=0.3]{\EUSocialFundLogoColorPictureEN}

    \end{textblock}





    \begin{textblock}{11}(1,4)

      \begin{flushleft}

        \mdseries

        \footnotesize

        \RaggedRight

        \color{jFrametitleFGColor}

        The content of this lecture is made available under a Creative
        Commons licence (\textsc{cc}), giving the author the credits
        (\textsc{by}) and putting an obligation to share on the same terms
        (\textsc{sa}). Figures and diagrams included in the lecture are
        authored by Paweł Węgrzyn et~al., and are available under the same
        license unless indicated otherwise. \\ The presentation uses the
        Beamer Jagiellonian theme based on Matthias Vogelgesang’s
        Metropolis theme, available under license \LaTeX{} Project
        Public License~1.3c at: \colorhref{https://github.com/matze/mtheme}
        {https://github.com/matze/mtheme}.

        Typographic design: Iwona Grabska-Gradzińska;
        \LaTeX{} Typesetting: Kamil Ziemian \\
        Proofreading: Leszek Hadasz, Wojciech Palacz,
        Monika Stawicka \\
        3D Models: Dariusz Frymus, Kamil Nowakowski \\
        Figures and charts: Kamil Ziemian, Paweł Węgrzyn,
        Wojciech Palacz \\
        Film editing: Agencja Filmowa Film \& Television Production~--
        Zbigniew Masklak \\
        Film editing: IMAVI -- Joanna Kozakiewicz, Krzysztof Magda, Nikodem
        Frodyma


      \end{flushleft}

    \end{textblock}

  \end{frame}





  \begin{frame}[standout]


    \begingroup

    \color{jFrametitleFGColor}

    #1

    \endgroup

  \end{frame}
}











% ------------------------------------------------------
% BibLaTeX
% ------------------------------------------------------
% Package biblatex, with biber as its backend, allow us to handle
% bibliography entries that use Unicode symbols outside ASCII.
\usepackage[
language=polish,
backend=biber,
style=alphabetic,
url=false,
eprint=true,
]{biblatex}

\addbibresource{Systemy-operacyjne-Bibliography.bib}





% ------------------------------------------------------
% Importing packages, libraries and setting their configuration
% ------------------------------------------------------





% ------------------------------------------------------
% Local packages
% ------------------------------------------------------
% Local configuration of this particular presentation
\usepackage{./Local-packages/local-settings}

\usepackage{./Local-packages/PGF-TikZ-Arrows-styles}

\usepackage{./Local-packages/PGF-TikZ-Diagram-styles}










% ####################################################################
\title{Systemy operacyjne}
\subtitle{Dod.~1. Z~historii systemów operacyjnych}

\author{Kamil Ziemian \\
  \email}


% \date{}
% ####################################################################










% ####################################################################
% Beginning of the document
\begin{document}
% ####################################################################





% ######################################
% Number of chars: 43k+, 81k+,
% Text is adjusted to the left and words are broken at the end of the line.
\RaggedRight
% ######################################





% ######################################
\maketitle
% ######################################





% ######################################
\begin{frame}
  \frametitle{Spis treści}


  \tableofcontents

\end{frame}
% ######################################










% ################################################
\section{Materiały do nauki}
% ################################################


% ######################################
\begin{frame}
  \frametitle{Do innej prezentacji}

  \textit{Wszystko to [omawiany program komputerowy] wydało mi~się
    aż~zbyt sprytne, a~źle zrozumiany spryt jest w~programowaniu bardzo
    złym pomysłem.} [...]

  \textit{Stało~się to podstawą naszej pierwszej reguły stylu
    programowania:} \\
  \textit{„Pisz przejrzyście~-- nie bądź za sprytny.} \\[0.3em]
  Str.~$29$, \parencite{Kernighan-Jak-Unix-tworzyl-historie-Pub-2021}.

\end{frame}
% ######################################





% ######################################
\begin{frame}
  \frametitle{UNIX i cała reszta}


  \textit{„BŁĘDY: polecenie} \texttt{rm} \textit{zapewne powinno pytać, czy
    plik tylko do odczytu, ma zostać usunięty.”} \\
  \hspace{5em} \textit{--~fragment instrukcji obsługi polecenia}
  \texttt{rm}\textit{, $3$~listopada $1971$~roku.} \\[0.3em]
  Str.~$61$, \parencite{Kernighan-Jak-Unix-tworzyl-historie-Pub-2021}.

  \textit{Joe Ossanna napisał program o~nazwie Nroff (ang. „new Roff”),
    w~nawiązaniu do~istniejącego programu do formatowania tekstu
    o~nazwie Roff.} \\[0.3em]
  Str.~$63$, \parencite{Kernighan-Jak-Unix-tworzyl-historie-Pub-2021}.

  \textit{Przy okazji warto wspomnieć, że~wywołanie systemowe}
  \texttt{creat} \textit{rzeczywiście jest zapisywane w~ten sposób~--
    wyłącznie z~powodu osobistych preferencji Kena Thompsona. Rob Pike
    spytał kiedyś Kena, co zmieniłby w~Uniksie, gdyby miał pisać go od~nowa.
    Co usłyszał? Napisałbym} \texttt{creat} \textit{z}~\texttt{e}
  \textit{na końcu.”}. \\[0.3em]
  Str.~$87$, \parencite{Kernighan-Jak-Unix-tworzyl-historie-Pub-2021}.

\end{frame}
% ######################################





% ######################################
\begin{frame}
  \frametitle{Do innej prezentacji}


  \textit{Język ten [Mowa o~języku C.] powstał na~początku lat~$70$.
    Opierał~się na doświadczeniach Dennisa [Ritchiego] z~implementacjami
    języków wysokiego poziomu dla Multicsa, lecz był znacznie mniej
    rozbudowany, ponieważ większość ówczesnych komputerów miała
    bardzo ograniczone możliwości; po prostu nie starczało im pamięci
    ani mocy obliczeniowej do obsługi kompilator równie skomplikowanego
    języka. Ten wymuszony minimalizm wpisywał~się w~preferencje
    Kena i~Dennisa dotyczące prostych, jednolitych mechanizmów. Język~C
    był też dobrze dopasowany do realnego sprzętu komputerowego~--
    było jasne, jak przetłumaczyć go na dobry, efektywnie działający
    kod.} \\[0.3em]
  Str.~$79$, \parencite{Kernighan-Jak-Unix-tworzyl-historie-Pub-2021}.

  Warto zaznaczyć, że~w~procesie narodzin języka~C, przełom~lat $60$-tych
  i~$70$-tych XX wieku, dużą rolę odegrał komputer PDP-$7$, który miał
  mniej niż $20$ kilobajtów całkowitej pamięci. Ilustruje to do pewnego
  stopnia skalę ograniczeń, które ukształtowały język~C.
  Por. str.~$50$--$51$, $78$--$80$
  \parencite{Kernighan-Jak-Unix-tworzyl-historie-Pub-2021}.

\end{frame}
% ######################################





% ######################################
\begin{frame}
  \frametitle{UNIX i cała reszta}


  \textit{Nazwa „grep” pochodzi od~dostępnego w~edytorze tekstowym}
  \texttt{ed} \textit{polecania} \texttt{g/re/p}, \textit{które
    wyświetla wszystkie linie pasujące do~wyrażenia regularnego}
  \texttt{re}\textit{;} \\[0.3em]
  Str.~$96$, \parencite{Kernighan-Jak-Unix-tworzyl-historie-Pub-2021}.

  \textit{Pewnym wariantem wyrażeń regularnych są symbole wieloznaczne
    dla nazw plików. Choć symbole wieloznaczne~są interpretowane
    przez powłokę, to ponieważ główna pamięć PDP-$7$ była bardzo
    ograniczone [Była mniejsza niż $20$ kilobajtów.], w~pierwszej
    implementacji odpowiadał za to osobny program o~nazwie}
  \texttt{glob} \textit{(od „globalny”) wywoływany przez powłokę,
    a~operację tworzenia listy plików na podstawie podanego wzorca
    nazywano „globbingiem”. Nazwa} \texttt{glob} \textit{do~dziś
    funkcjonuje w~bibliotekach kilku języków programowania, między
    innymi w~Pythonie.} \\[0.3em]
  Str.~$99$, \parencite{Kernighan-Jak-Unix-tworzyl-historie-Pub-2021}.

\end{frame}
% ######################################





% ######################################
\begin{frame}
  \frametitle{UNIX i cała reszta}


  \textit{Jedną z~pierwszych zasług Ala (Alfreda) Aho jeśli chodzi
    o~rozwój Uniksa było napisanie rozszerzenia do grepa, które pozwalało
    na~stosowanie bogatszych wyrażeń regularnych, przykładowo pozwalających
    na wyszukiwanie alternatyw w~rodzaju} \texttt{to|tamto}\textit{.}
  \textit{Al nazwał swój program egrep od} extended
  grep\textit{.} \\[0.3em]
  Str.~$100$, \parencite{Kernighan-Jak-Unix-tworzyl-historie-Pub-2021}.

  \textit{BCPL (ang. } Basic Combined Programming language\textit{)
    był kolejnym językiem z~założenia przeznaczonym do~programowania
    systemów. Został opracowany przez Martina Richardsa, profesora
    z~Uniwersytetu w~Cambridge, który podczas gościnnej pracy w~M.I.T.
    w~$1966$ roku napisał dla niego kompilator. BCPL był znacznie
    prostszy od~dowolnego dialektu PL/I i~dobrze przystosowany do~pisania
    kodu systemu operacyjnego. Członkowie zespołu Bell Labs pracujący
    nad Multicsem byli znakomicie zaznajomieni z~BCPL.} \\[0.3em]
  Str.~$101$, \parencite{Kernighan-Jak-Unix-tworzyl-historie-Pub-2021}.

\end{frame}
% ######################################





% ######################################
\begin{frame}
  \frametitle{UNIX i cała reszta}


  \textit{Gdy laboratorium Bell wycofało~się z~prac nad Multicsem
    Ken Thompson uznał, że~„żaden komputer nie może~się objeść bez
    Fortrana” i~zaczął pisać kompilator Fortrana dla PDP-$7$.
    Wyzwanie okazało się trudne, ponieważ Unix na PDP-$7$ miał
    tylko $4\text{K}$ $18$-bitowych słów ($8$ kilobajtów) pamięci
    podstawowej na~programy użytkownika, takie jak kompilator.}

  \textit{Ken nie poddawał~się i~w~wyniku nieustanych modyfikacji otrzymał
    język, który mieścił~się w~PDP-$7$ i~był bardzie podobny do BCPL niż
    do~Fortrana. Nazwał go~B. Jak wyjaśnił Dennis Ritchie w~publikacji
    „The~Development~of the~C Language” w~$1993$ roku}: \\
  \textit{„B~można porównać do~C bez typów;}
  \\[0.3em]
  Str.~$101$, \parencite{Kernighan-Jak-Unix-tworzyl-historie-Pub-2021}.

\end{frame}
% ######################################





% % ######################################
% \begin{frame}
%   \frametitle{Materiały do nauki}



% \end{frame}
% % ######################################





% % ######################################
% \begin{frame}
%   \frametitle{Materiały do nauki}



% \end{frame}
% % ######################################






% % ######################################
% \begin{frame}
%   \frametitle{Materiały do nauki}



% \end{frame}
% % ######################################





% % ######################################
% \begin{frame}
%   \frametitle{Czy ten przedmiot będzie trudny?}



% \end{frame}
% % ######################################





% % ######################################
% \begin{frame}
%   \frametitle{Włączanie komputera, poziom normalnego użytkownika}




% \end{frame}
% % ######################################





% % ######################################
% \begin{frame}
%   \frametitle{Kilka pojęcia}



% \end{frame}
% % ######################################





% % ######################################
% \begin{frame}
%   \frametitle{Kilka pojęcia}




% \end{frame}
% % ######################################





% % ######################################
% \begin{frame}
%   \frametitle{Kilka pojęcia}




% \end{frame}
% % ######################################





% % ######################################
% \begin{frame}
%   \frametitle{Uruchamianie komputera z~systemem Pentium}




% \end{frame}
% % ######################################





% % ######################################
% \begin{frame}
%   \frametitle{Uruchamianie komputera z~systemem Pentium}


%   \begin{itemize}

%   \item[5)] \textsc{bios} testuje po kolei urządzenia rozruchowe
%     z~wspomnianej wcześniej listy, aż~znajdzie pierwsze, który zawiera
%     działający system operacyjny.

%   \item[6)] \textsc{bios} wczytuje pierwszy sektor ze~znalezionego
%     w~poprzednim punkcie urządzenia rozruchowego do pamięci i~go uruchamia.

%   \item[7)] Program z~pierwszego sektora sprawdza zapisaną na jego końcu
%     listę partycji, by~ustalić która z~nich jest partycją aktywną.
%     Następnie wczytuje z~tej partycji pomocniczy program rozruchowy.

%   \item[8)] Pomocniczy program rozruchowy wczytuje system operacyjny
%     z~aktywnej partycji i~go uruchamia.

%   \item[9)] System operacyjny odczytuje informacje konfiguracyjne z~systemu
%     \textsc{bios}. Dla każdego dostępnego urządzenia sprawdza, czy posiada
%     do niego sterowniki. Jeśli nie, to prosi o~ich zainstalowanie
%     z~odpowiedniego źródła.

%   \end{itemize}

% \end{frame}
% % ######################################





% % ######################################
% \begin{frame}
%   \frametitle{Uruchamianie komputera z~systemem Pentium}


%   \begin{itemize}

%   \item[10)] Jeśli system operacyjny dysponuje wszystkimi sterownikami,
%     to ładuje je do jądra systemu.

%   \item[11)] System operacyjny tworzy tabele systemowe oraz procesy
%     działające w~tle.

%   \item[12)] Uruchamiane jest okno logowania.

%   \end{itemize}

% \end{frame}
% % ######################################






% % ######################################
% \begin{frame}
%   \frametitle{Bootowanie}


%   W~literaturze funkcjonuje termin \textbf{bootwoanie}, zwane też
%   \textbf{uruchamianiem} lub \textbf{rozruchem}. Odnosi~się ono albo do
%   całej procedury uruchamiania komputer opisanej powyżej, albo tylko
%   stawiania systemu operacyjnego, czyli procedury przeprowadzonej
%   od momentu, gdy \textsc{bios} wczytał pierwszy jego sektor do pamięci
%   (punkt siedem i~dalej). Acz to pojęcie nie jest specjalnie ostro
%   zdefiniowane.

% \end{frame}
% % ######################################





% % ######################################
% \begin{frame}
%   \frametitle{Czy uruchomienie komputera jest proste czy trudne?}


%   Zależy jak do tego podchodzimy. I~tak jest z~większością rzeczy
%   w~informatyce.

% \end{frame}
% % ######################################










% % ################################################
% \section{Zdalne logowanie~się do GNU/Linuxa z~Windowsa}
% % ################################################


% % ######################################
% \begin{frame}
%   \frametitle{Korzystanie z~systemu Artemis}


%   Wiemy, że~dla wielu osób zainstalowanie na swoim komputerze systemu
%   GNU/Linux, może być sporym wyzwaniem, dlatego tutaj wyjaśnimy, jak
%   można~się zalogować do niego z~poziomu systemu Windows. Dokładniej,
%   pokażemy, jak połączyć~się przez internet do komputerów dysponujących
%   systemem GNU/Linux, poprzez uczelniany system Artemis. Chciałbym przy tym
%   bardzo podziękować panu Tomaszowi Molęckiemu z~działu \textsc{it}
%   \textsc{wsz}i\textsc{b}u za~udzielenie nam informacji na temat
%   procedury logowania~się do tego systemu.

%   Pierwszy krok polega na pobraniu programu
%   \colorhref{https://winscp.net/eng/download.php}{Win\textsc{scp}}. Gdy
%   już mamy go dostępnego, znajdujemy okno logowania, które wypełniamy
%   zgodnie ze~zdjęciem widocznym na następnym slajdzie. Dane logowania
%   wpisujemy takie same jak dla systemu SUSZI. Gdy to jest gotowe, klikamy
%   na przycisk logowania. W~razie problemów proszę pisać do nas (\email),
%   lub pana Molęckiego.

% \end{frame}
% % ######################################





% % ######################################
% \begin{frame}
%   \frametitle{Korzystanie z~systemu Artemis}

%   \vspace{-0.5em}


%   \begin{figure}

%     \centering


%     \includegraphics[scale=0.45]
%     {./Presentations-pictures/Miscancellous-pictures/Artemis.jpg}


%     \caption{Logowanie~się do systemu Artemis za pomocą programu
%       Win\textsc{scp}.}


%     \label{fig:Artemis-system}

%   \end{figure}

% \end{frame}
% % ######################################












% % ################################################
% \section{Włączanie powłoki BASH}
% % ################################################


% % ######################################
% \begin{frame}
%   \frametitle{Powłoka BASH}

%   \vspace{-0.5em}


%   \begin{figure}

%     \centering


%     \includegraphics[scale=0.23]
%     {./Presentations-pictures/Miscancellous-pictures/BASH-shell.png}


%     \caption{Przykładowy wygląd włączonej powłoki \textsc{bash}, naszego
%       głównego narzędzia pracy na tym przedmiocie.}


%     \label{fig:BASH-shell}

%   \end{figure}

% \end{frame}
% % ######################################





% % ######################################
% \begin{frame}
%   \frametitle{Postęp technologiczny}

%   \vspace{-0.5em}


%   \begin{figure}

%     \label{fig:Evolution-of-OS}

%     \centering


%     \includegraphics[scale=0.3]
%     {./Presentations-pictures/Miscancellous-pictures/Evolution-of-operating-systems.jpg}

%   \end{figure}

% \end{frame}
% % ######################################





% % ######################################
% \begin{frame}
%   \frametitle{Włączanie powłoki}


%   \textit{Gdy samochody rozwijały~się tak jak komputery, to dziś
%     Rolls-Royce kosztowałby jakieś 100\$, przejeżdżałaby milion mil na
%     jednym galonie benzyny i~raz do roku eksplodował, zabijając wszystkich
%     w~środku.} \\
%   Robert Cringely, \textit{InfoWorld}, tłum. swobodne. Por.~str.~221
%   \parencite{Garfinkel-Weise-Strassmann-The-UNIX-HATERS-Handbook-Pub-1994}.

%   Naszym główny narzędziem na tych zajęciach będzie powłoka
%   (ang. \textit{shell}) \textsc{bash}, więc upewnijmy~się, że~umiemy ją
%   włączyć, co na samy początku może być naprawdę niewdzięcznym zadaniem.
%   W~pierwszym kroku proszę wcisnąć \texttt{Ctrl-Alt-t} i~zobaczyć, czy~się
%   otworzy.

%   Jeśli to nie zadziała, to proszę pochodzić kursorem po ekranie,
%   aż~znajdziemy okienko „Szukaj”, „Wyszukaj” lub jeszcze jakaś inna nazwa.
%   Proszę tam wpisać „Terminal”, „Konsola” lub „Consol” i~zobaczyć, czy
%   wyświetli~się odpowiednia ikona do kliknięcia.

%   % !!!!!!! Nie usuwać.
%   % Istnieje wiele wartych uwagi materiałów do nauki \textsc{bash}a, niekiedy
%   % dość już wiekowych. Przykładowo, można sięgnąć do tutorialu
%   % Gilesa Orra
%   % \colorhref{https://www.gilesorr.com/bashprompt/howto/}
%   % {\textit{Bash Prompt HOWTO}} \parencite{Orr-Bash-Prompt-HOWTO-Ver-2009}.

% \end{frame}
% % ######################################










% % ##################
% \begin{frame}
%   \frametitle{Z~BASHem dzieje~się coś dziwnego}


%   Błędne wprowadzenie jakiegoś polecenia może spowodować, że~\textsc{bash}
%   przestanie reagować, lub będzie reagował zupełnie nie tak jak tego
%   oczekujemy. Ponieważ \textsc{bash} nie wiem czym jest litość, w~zależności
%   od tego co konkretnie~się zepsuło, dokładna procedura naprawy jest inna.

%   Jednak w~większości przypadków, by wrócić do normalnej pracy wystarczy
%   wcisnąć \alert{raz} \texttt{Ctrl-d}. Należy zaznaczyć, że~wciśnięcie
%   \texttt{Ctrl-d} \alert{więcej niż raz} może spowodować zamknięcie powłoki
%   \textsc{bash} i~potrzeby zrobienia kilku rzeczy ponownie. Z~tego też
%   powodu, początkujący \alert{nigdy} nie powinni używać tego skrótu, jeśli
%   wszystko działa jak trzeba.

%   Mimo tych zastrzeżeń, \texttt{Ctrl-d} jest jednym z~naszych najlepszych
%   przyjaciół.

% \end{frame}
% % ##################










% % ######################################
% \section{Podstawy pracy z~poziomu powłoki BASH}
% % ######################################



% % ##################
% \begin{frame}
%   \frametitle{Kilka słów wstępnych i~ostrzeżeń}


%   Przed chwilą mówiliśmy, że~\textsc{bash} nie zna słowa litość i~każda
%   drobna zmiana może wszystko zniszczyć. Jednocześnie nasza swoboda nie jest
%   całkiem zerowa (choć może często sprawiać takie wrażenie), więc
%   podając tutaj różne polecenia, musimy zdecydować~się na pewno konwencje.
%   Proszę więc mieć świadomość, że~\alert{pewne} zmiany można wprowadzić
%   do~poleceń i~nie zmienią one sposób ich działania.

%   Jednak zasady które mówią co \alert{można} zmienić w~poleceniu
%   \textsc{bash}a, bez obrócenia wszystkiego w~ruiny, są często nieintuicyjne
%   i~skomplikowane. Ponieważ zaś zawsze $1000$ różnych rzeczy może zepsuć
%   dane polecenie, zdecydowaliśmy, że~w~tym kursie będziemy starali~się
%   podawać te wersje poleceń, które my uważamy za~optymalne i~które działają
%   poprawnie, nie zaś wnikać w~pokrętne zasady „możesz zmienić to tak i~tak,
%   ale~już nie w~ten sposób”.

% \end{frame}
% % ##################





% % ##################
% \begin{frame}
%   \frametitle{Kilka słów wstępnych i~ostrzeżeń}


%   Testowanie różnych wariantów tych poleceń, sprawdzanie który ich wariant
%   działa, a~który nie, jest oczywiście czymś dobrym, co pozwala lepiej
%   zrozumieć powłokę \textsc{bash}, więc bardzo nasz cieszy, gdy ktoś tego
%   próbuje. Zajmuje to jednak czas, a~większość ludzi pragnie by „komputer
%   zrobił to co trzeba”, nie zaś zgłębiać wszystkie niuanse jego pracy.
%   Biorąc to wszystko pod uwagę, uznaliśmy, że~obecne sposób wykładu jest
%   dobrym wyborem.

%   Wszystko utrudnia dodatkowy fakt, że~od czasu do czasu można~się natknąć
%   na~dwa warianty poleceń \textsc{bash}a, które działają tak samo, tylko,
%   że~to nie do końca prawda. Działają tak samo w~$90\%$ przypadków,
%   w~pozostały~$10\%$, działają inaczej. By uniknąć niekończących~się
%   dywagacji na temat dzielenia włosa na czworo w~\textsc{bash}u, proszę
%   przyjąć, że~jeśli polecenie nie działa jak trzeba, należy~się upewnić,
%   iż~są one \alert{dokładnie} takie jak podano tutaj.

% \end{frame}
% % ##################





% % ##################
% \begin{frame}
%   \frametitle{Kilka słów wstępnych i~ostrzeżeń}


%   Do problemu, której wersji poleceń/kodu \textsc{bash}a należy używać,
%   wrócimy jeszcze przy okazji pisania skryptów dla tej powłoki. To jednak
%   na tą chwilę dość odległy temat.

%   Pomimo tego, że~staramy~się dokładnie przetestować wszystkie podane
%   w~tych wykładach polecenia \textsc{bash}a, nie możemy wykluczyć tego,
%   że~popełnimy jakiś błąd. Dlatego bardzo prosimy by każdy, kto taki
%   znajdzie, przesłał go nam pod adres \email.

% \end{frame}
% % ##################





% % ##################
% \begin{frame}
%   \frametitle{Zaczynamy}


%   Zacznijmy od przyjęcia, że~pracujemy na komputerze jako użytkownik,
%   którego nazwa w~systemie operacyjnym została wybrana jako \texttt{adam}.
%   Jeśli włączymy konsolę zobaczymy pewnie coś takiego: \\
%   \texttt{adam:$\sim$\$} \\
%   Od tego momentu będziemy zwykle skracali ten znak zachęty \\
%   (ang.~\textit{prompt}) do symbolu \texttt{\$}. Przepisywanie pełnej jego
%   postaci wnosi bardzo mało informacji, za~to zabiera sporo miejsca.

%   Na początku zajmiemy~się tym, jak poruszać~się po systemie plików
%   i~katalogów. Na pewno wiedzą Państwo, że~na komputerze znajdują~się
%   katalogi (ang.~\textit{directories}), zwane też folderami, które zawierają
%   pliki i~inne katalogi. Z~punktu widzenia systemu GNU/Linux
%   katalog jest plikiem, jedynie specjalnego rodzaju. Dla tego systemu
%   wszystko jest plikiem lub procesem.

% \end{frame}
% % ##################





% % ##################
% \begin{frame}
%   \frametitle{Polecenia}


%   Podstawy działania katalogów w~systemie GNU/Linux są w~zasadzie takie
%   same jak w~systemie Windows, musimy~się jedynie przyzwyczaić do innych
%   konwencji używanych w~tym systemie. Dużo większym wyzwaniem zapewne
%   będzie zamienienie środowiska graficznego (\textsc{gui},
%   ang.~\textit{Graphic User Interface}) na pracę z~poziomu powłoki.

%   Podstawą pracy w~powłoce \textsc{bash} jest korzystanie z~poleceń (ang.
%   \textit{commands}). Bez wnikania w~szczegóły, każde polecenie posiada
%   nazwę i~wykonuje pewne działanie. By wykonać dane polecenie w~powłoce
%   należy najpierw wpisać jej nazwę, następnie zaś jej odpowiednie
%   opcje (flagi) i~argumenty. Część poleceń można wykonać bez wpisywania
%   opcji czy~argumentów.

%   Trzeba~się niestety przyzwyczaić do tego, że~nazwy poleceń, jak i~nazwy
%   ich opcji tych poleceń w~systemie GNU/Linux potrafią być dość dziwne
%   i~trudne do~zapamiętania.

% \end{frame}
% % ##################





% % ##################
% \begin{frame}
%   \frametitle{Polecenia}


%   Podstawowych poleceń systemu GNU/Linux jest około~$100$, do tego każde ma
%   kilkanaście albo więcej opcji. Z~tego powodu nauka ich jest bardzo
%   czasochłonna i~niełatwa. Teoretycznie w~każdej chwili w~\textsc{bash}u
%   możemy skorzystać z~manuala (podręcznika), gdzie są opisy wszystkich
%   poleceń. W~praktyce, żeby zrozumieć jego treść, wymagana jest już dość
%   spora wiedza.

%   Zdecydowaliśmy~się przyjąć następujące podejście. Naszym głównym celem
%   jest przekazanie Państwu umiejętności pracy spod powłoki \textsc{bash}
%   w~możliwie krótki i~bezbolesny sposób. Dlatego też, dość szybko
%   wprowadzimy możliwie dużo poleceń, bez zagłębiania~się w~zbyt wiele
%   szczegółów. Niekiedy będziemy wracać do wcześniej już omówionych poleceń,
%   by uzupełnić wiadomości o~nich. Wyczerpujące omówienie wszystkich
%   podstawowych poleceń i~tak zabrałoby zbyt dużo czasu, to zaś podejście
%   uważamy, za rozsądny kompromis.

% \end{frame}
% % ##################





% % ##################
% \begin{frame}
%   \frametitle{Polecenie \texttt{ls}}


%   Zacznijmy od~polecenia
%   \colorhref{https://linux.fandom.com/pl/wiki/Ls}{\texttt{ls}}. Jego nazwa
%   pochodzi prawdopodobnie od~ang.~\textit{List fileS}, ale są różne opinie
%   na ten temat. Polecenie to wyświetla zawartość katalogu w~którym~się
%   znajdujemy, z~tego powodu wynik działania tego polecenia będzie~się cały
%   czas zmieniać, jednym z~możliwych może być ten przedstawiony poniżej. \\
%   \texttt{\$ ls} \\
%   \texttt{{\color{jAxisBlue} Dokumenty}} \quad
%   \texttt{introduction-to-GNU-linux.pdf} \quad
%   \texttt{{\color{jAxisBlue} Pobrane}}

%   Polecenie \texttt{ls} jest przykładem takiego polecenia, które można
%   uruchomić bez dodatkowych opcji czy argumentów, ono samo ustali w~jakim
%   katalogu jesteśmy i~wypisze jego zawartość.

%   Obecnie jest standardem, że~powłoka używa kolorów i~różnych typów fontów,
%   by~informować nas o~typie i~własnościach pliku. W~przykładzie powyżej
%   nazwy katalogów są wyróżnione na niebiesko, trzeba jednak pamiętać,
%   że~wszystko to~zależy od~konkretnych ustawień systemu operacyjnego
%   i~łatwo może zostać zmienione.

% \end{frame}
% % ##################





% % ##################
% \begin{frame}
%   \frametitle{Polecenie \texttt{pwd}}


%   Kolejnym poleceniem jest
%   \colorhref{https://www.arturpyszczuk.pl/commands-pwd.html}{\texttt{pwd}}
%   od ang.~\textit{Print Working Directory}, rezultat jego działania można
%   zobaczyć poniżej. \\
%   \texttt{\$ pwd} \\
%   \texttt{/home/adam}

%   Ta krótka linia zawiera bardzo dużo informacji i~potrzebujemy przynajmniej
%   kilka z~nich wyjaśnić. Sposób w~jaki pliki i~katalogi są ułożone na~dysku
%   jest zdefiniowane przez \textbf{system plików} danego systemu
%   operacyjnego i~ten używany w~systemie GNU/Linux różni~się mocno tego
%   w~systemie Windows.

%   Zacznijmy od kilku uwag odnośnie notacji. Jeśli pracowalibyśmy w~tym
%   ostatnim systemie i~na dysku~C znajdowałby~się
%   katalog \texttt{Systemy-operacyjne-Ćwiczenia} a~w~nim katalog
%   \texttt{Zajęcia-01}, to ścieżka do niego miałaby postać: \\
%   \texttt{C:\textbackslash\textbackslash
%     Systemy-operacyjne-Ćwiczenia\textbackslash Zajęcia-01} \\

% \end{frame}
% % ##################





% % ##################
% \begin{frame}
%   \frametitle{System plików}


%   Odpowiednikiem ścieżki z~systemu Windows: \\
%   \texttt{C:\textbackslash\textbackslash
%     Systemy-operacyjne-Ćwiczenia\textbackslash Zajęcia-01} \\
%   w~systemie GNU/Linux byłaby ścieżka: \\
%   \texttt{/home/adam/Systemy-operacyjne-Ćwiczenia/Zajęcia-01} \\
%   Widzimy więc, że~gdy Windows do oddzielania nazw katalogów i~plików
%   używa symbolu „backslasha”: „\texttt{\textbackslash}”, to GNU/Linux
%   używa symbolu „slasha”: „\texttt{/}”. Jest to stosunkowo mała różnica,
%   ale trzeba~się do niej przyzwyczaić.

%   Poważniejsza różnica jest następująca. Komputery z~systemem GNU/Linux
%   mogą posiadać wiele dysków, nie jest to żadnym problemem, ale~system
%   plików GNU/Linuxa nie zna pojęcia dysku jako takiego. W~systemie GNU/Linux
%   istnienie naczelny katalog zwany rootem, czyli po polsku „korzeniem”,
%   w~którym zawierając~się wszystkie inne pliki oraz katalogi, nieważne na
%   którym fizycznym dysku są one umieszczone. Katalog root oznaczamy
%   symbolem~„\texttt{/}”.

% \end{frame}
% % ##################





% % % ##################
% % \begin{frame}
% %   \frametitle{System plików}


% %   Odpowiednikiem ścieżki z~systemu Windows: \\
% %   \texttt{C:\textbackslash\textbackslash
% %     Systemy-operacyjne-Ćwiczenia\textbackslash Zajęcia-01} \\
% %   w~systemie GNU/Linux byłaby ścieżka: \\
% %   \texttt{/home/adam/Systemy-operacyjne-Ćwiczenia/Zajęcia-01} \\
% %   Widzimy więc, że~gdy Windows do oddzielania nazw katalogów i~plików
% %   używa symbolu „backslasha”: „\texttt{\textbackslash}”, to GNU/Linux
% %   używa symbolu „slasha”: „\texttt{/}”. Jest to stosunkowo mała różnica,
% %   ale trzeba~się do niej przyzwyczaić.

% %   Poważniejsza różnica jest następująca. Komputery z~systemem GNU/Linux
% %   mogą posiadać wiele dysków, nie jest to żadnym problemem, ale~system
% %   plików GNU/Linuxa nie zna pojęcia dysku jako takiego. W~systemie GNU/Linux
% %   istnienie naczelny katalog zwany rootem, czyli po polsku „korzeniem”,
% %   w~którym zawierając~się wszystkie inne pliki oraz katalogi, nieważne na
% %   którym fizycznym dysku są one umieszczone. Katalog root oznaczamy
% %   symbolem~„\texttt{/}”.

% % \end{frame}
% % % ##################





% % ##################
% \begin{frame}
%   \frametitle{System plików}


%   Wróćmy do polecenia \\
%   \texttt{\$ pwd} \\
%   \texttt{/home/adam} \\
%   Linię \texttt{/home/adam} czytamy „W~katalogu root („\texttt{/}”)
%   znajduje~się katalog \texttt{home}, w~nim zaś znajduje~się katalog
%   \texttt{adam}.” Polecenie \texttt{pwd} podaje więc w~którym katalogu~się
%   obecnie znajdujemy, za~pomocą ścieżki, która zaczyna~się od katalogu root:
%   „\texttt{/}”.

%   Katalog w~którym obecnie jesteśmy nosi nazwę \textbf{katalogu roboczego}
%   (ang.~\textit{working directory}), a~jak powiedzieliśmy \texttt{pwd} jest
%   akronimem od~\textit{Print Working Directory}, czyli „wyświetl katalog
%   roboczy”.

%   Ścieżkę podającą lokalizację katalogu lub pliku nazywamy
%   \textbf{ścieżką bezwzględną}, jeśli zaczyna~się ona od katalogu root.
%   W~przeciwnym razie mówimy o~\textbf{ścieżce względnej} (podawanej
%   względem innego katalogu niż root).

% \end{frame}
% % ##################





% % ##################
% \begin{frame}
%   \frametitle{System plików}


%   Bez wielkich uproszczeń można przyjąć, że~\textsc{bash}u \alert{każdy}
%   ciąg symboli zaczynający od~symbolu „\texttt{/}” zostanie zinterpretowany
%   jako ścieżka do jakiegoś pliku lub katalogu. Z~tego powodu trzeba
%   \alert{bardzo} uważać przy stosowaniu symbolu~„\texttt{/}”.

%   W~szczególności, gdy stosujemy symbol „\texttt{/}”, trzeba zwracać dużą
%   uwagę na~spacje. Ścieżka postaci \\
%   \texttt{/home/adam} \\
%   jest poprawna. Ale już ścieżka \\
%   \texttt{/ home/adam} \\
%   jest \alert{błędna}! To w~jaki sposób jest błędna, to omówimy innym
%   razem. Doświadczenie uczy, że~na początku sprawia to ludziom
%   bardzo dużo problemów.

%   To, że~główny katalog nosi nazwę „korzenia” (ang.~\textit{root}) nie jest
%   przypadkiem, ale do tej kwestii wrócimy innym razem.

% \end{frame}
% % ##################





% % ##################
% \begin{frame}
%   \frametitle{Katalog domowy}


%   W~systemie GNU/Linux każdy użytkownik ma swój katalog domowy, w~którym
%   może przechowywać potrzebne mu pliki. Standardowo, choć mogą być wyjątki,
%   użytkownik \texttt{adam} otrzyma katalog \texttt{/home/adam}, użytkownik
%   \texttt{barbara} katalog \texttt{/home/barbara},~etc. Katalog domowy
%   bieżącego użytkownika posiada też alternatywną nazwę „$\sim$”. Stąd znak
%   zachęty \\
%   \texttt{adam:$\sim$\$} \\
%   przekazuje nam informację, że~działamy jako użytkownik \texttt{adam}
%   i~znajdujemy~się w~jego katalogu domowym.

%   Załóżmy, że~katalogu domowym znajduje~się katalog o~nazwie
%   \texttt{Systemy-operacyjne-Ćwiczenia}. By sprawdzić, czy tak jest
%   naprawdę możemy użyć polecenia \texttt{ls}: \\
%   \texttt{\$ ls} \\
%   \texttt{{\color{jAxisBlue} Dokumenty}} \quad
%   \texttt{introduction-to-GNU-linux.pdf} \quad
%   \texttt{{\color{jAxisBlue} Pobrane}} \\
%   \texttt{{\color{jAxisBlue} Systemy-operacyjne-Ćwiczenia}} \\

% \end{frame}
% % ##################





% % ##################
% \begin{frame}
%   \frametitle{Polecenie \texttt{cd}}


%   Jeśli w~bieżącym katalogu znajduje~się katalog
%   \texttt{Systemy-operacyjne-Ćwiczenia} to możemy do niego wejść
%   używając polecenia
%   \colorhref{https://linux.fandom.com/pl/wiki/Cd}{\texttt{cd}}
%   (od ang.~\textit{Change Directory}): \\
%   \texttt{adam:$\sim$\$ cd Systemy-operacyjne-Ćwiczenia} \\
%   \texttt{adam:$\sim$/Systemy-operacyjne-Ćwiczenia\$} \\
%   Wyjątkowo zapisaliśmy pełniejszą wersję znaku zachęty, by podkreślić,
%   iż~informuje on~nas w~którym katalogu jesteśmy. Mianowicie ścieżka
%   \texttt{$\sim$/Systemy-operacyjne-Ćwiczenia} informuje nas, że~w~katalogu
%   domowych użytkownika \texttt{adam} znajduje~się katalog
%   \texttt{Systemy-operacyjne-Ćwiczenia} i~w~nim właśnie jesteśmy.

%   Użyte wyżej polecenie \texttt{cd} jest przykładem polecenia, które
%   pobiera argument, mianowicie nazwę katalogu do którego możemy wejść.
%   W~istocie polecenie to pozwala nam na bardzo swobodne poruszanie~się
%   po systemie plików, jeśli opanujemy dobrze zapis ścieżek w~GNU/Linuxie,
%   ale wszystko po kolei.

% \end{frame}
% % ##################





% % ##################
% \begin{frame}
%   \frametitle{Polecenie \texttt{cd}}


%   Załóżmy, że~teraz chcielibyśmy przejść do katalogu nad nami. By to zrobić
%   należy przekazać poleceniu \texttt{cd} argument „\texttt{..}”: \\
%   \texttt{adam:$\sim$/Systemy-operacyjne-Ćwiczenia\$ cd ..} \\
%   \texttt{adam:$\sim$\$} \\
%   Dlaczego jako argument przesyłamy akurat „\texttt{..}” wyjaśnimy w~swoim
%   czasie, na razie przyjmijmy, iż~tak po prostu jest. Dla lepszej
%   ilustracji tego co~zrobiliśmy, ponownie wypisaliśmy znak zachęty,

%   Załóżmy, że~chcemy teraz wejść do katalogu
%   \texttt{Systemy-operacyjne-Ćwiczenia} i~utworzyć tam plik
%   \texttt{Notatki-z-zajęć.txt}. By utworzyć ten plik możemy skorzystać
%   z~polecenia
%   \colorhref{https://linux.fandom.com/pl/wiki/Touch}{\texttt{touch}}. Pełna
%   sekwencja poleceń będzie wyglądać tak. \\
%   \texttt{\$ cd Systemy-operacyjne-Ćwiczenia} \\
%   \texttt{\$ touch Notatki-z-zajęć.txt} \\

% \end{frame}
% % ##################





% % ##################
% \begin{frame}
%   \frametitle{Edycja plików z~pomocą programu gedit}


%   Możemy teraz użyć polecenia \texttt{ls} by~się przekonać, czy~plik
%   \texttt{Notatki-z-zajęć.txt} został poprawnie utworzony.
%   Z~komputerami nigdy nic nie wiadomo, więc warto to sprawdzić.

%   W~następnym kroku otworzymy go do edycji za~pomocą programu
%   \colorhref{https://gedit-text-editor.org/}{gedit}. By to zrobić
%   wpisujemy w~powłoce \\
%   \texttt{\$ gedit Notatki-z-zajęć.txt \&} \\
%   Proszę pamiętać o~dodaniu na końcu polecenia znaku
%   \alert{ampersand}:~„\&”. Inaczej gedit zablokuje nam możliwość wpisywania
%   kolejnych poleceń do powłoki i~by ją odzyskać, będziemy musieli
%   wyłączyć gedita albo użyć kombinacji poleceń \texttt{jobs} i~\texttt{bg}.

%   Czemu uruchomienie gedita bez symbolu ampersanda ma takie skutki, omówimy
%   gdy będziemy dyskutowali zagadnienie procesów i~tego jak działają
%   w~systemie GNU/Linux. Nie jest to najprostszy temat na świecie.

% \end{frame}
% % ##################





% % ##################
% \begin{frame}
%   \frametitle{Okienko gedita}


%   Jeśli wszystko poszło dobrze, to na ekranie zobaczymy okienko Gedita,
%   które wygląda mniej więcej tak jak to na~zdjęciu~\eqref{fig:Gedit-window}.





%   \begin{figure}

%     \centering

%     \includegraphics[scale=0.145]
%     {./Presentations-pictures/Miscancellous-pictures/gedit-window.png}

%     \caption{Otwarte okienko programu gedit.}


%     \label{fig:Gedit-window}

%   \end{figure}

% \end{frame}
% % ##################





% % ##################
% \begin{frame}
%   \frametitle{Edytory tekstów}


%   Teraz możemy do pliku \texttt{Notatki-z-zajęć.txt} wpisać tekst, choćby
%   \textit{Notatki z~przedmiotu „Systemy operacyjne”}. Proszę pamiętać, by
%   zapisywać wprowadzone zmiany, służy do tego choćby skrót~\texttt{Ctrl-s}.

%   Kilka słów dlaczego wybraliśmy do~edycji plików program gedit. Na~tym
%   kursie korzystamy głównie z~niego, gdyż jest on prosty w~obsłudze
%   i~w~kilka sekund po jego otwarciu użytkownik może już pisać odpowiedni
%   tekst. Zostawiamy jednak Państwu pełną swobodę w~wyborze edytora tekstu.
%   Nie jest ważne jakiego programu~się używa, ważne jest, żeby udało~się
%   wykonać zadanie jakie przed nami stoi.

%   Jeśli chodzi o~edytory dostępne pod systemem GNU/Linux, to w~naszej
%   ocenie warto zwrócić uwagę choćby na~dwa owiane już legendą programy
%   do edycji tekstu \colorhref{https://www.vim.org/}{vim}
%   i~\colorhref{https://www.gnu.org/software/emacs/}{GNU Emacs}.

% \end{frame}
% % ##################










% % ################################################
% \section{Czy tylko mi~się wydaje, że~ten system jest dziwny?}
% % ################################################


% % ######################################
% \begin{frame}
%   \frametitle{Czy tylko mi~się wydaje, że~ten system jest dziwny?}


%   Niejedna z~poznanych przez nas cech systemu operacyjnego, może
%   wydawać~się nam dziwna. Czy życie nie mogłoby być prostsze, a~system
%   operacyjny bardziej zrozumiały? Odpowiedź na oba pytania brzmi
%   „Prawdopodobnie tak”. Dlaczego więc tak nie jest?

%   Warto sobie uświadomić, że~niektórzy jeżdżą jeszcze zabytkowymi
%   samochodami wyprodukowanymi w~latach $60$-tych XX~wieku. Prawdopodobnie
%   gdzieś na świecie działają jakieś komputery wyprodukowane w~tych latach,
%   ale większa jest szansa zobaczenia go w~muzeum, niż wykonującego jakąś
%   pracę. W~przeciwieństwie do tego, pomysł z~lat $60$-tych mają~się
%   świetnie, również w~informatyce. Pracując z~systemami często
%   napotykamy~się na jakiś jego fragment, który wymyślono w~owym czasie,
%   odbija on więc ówczesny stan wiedzy i~możliwości techniczne. By więc
%   zrozumieć pewne osobliwości systemów operacyjnych warto udać~się w~krótką
%   podróż w~przeszłość.

% \end{frame}
% % ######################################





% % ######################################
% \begin{frame}
%   \frametitle{Trochę historii}


%   Główne przesłanie tej historii jest następujące. Systemy operacyjne
%   których dziś używamy, wyewoluowały do swojej obecnej formy formy
%   na przestrzeni dekad, wchłaniając i~łącząc pomysły pochodzące
%   z~różnorakich źródeł. Stąd ich dzisiejszy czas może sprawiać wrażenie
%   przestarzałego, bałaganiarskiego i~pełnego niekonsekwencji. Jest to prawdą
%   w~przypadku w~zasadzie wszystkich popularnych systemów operacyjnych, nie
%   tylko omawianego przez nas GNU/Linuxa.

%   Zaczynamy w~późnej epoce brązu komputerów, czyli roku~$1965$.
%   W~owym roku w~USA, uczelnia M.I.T oraz firmy General Electrics i~Bell
%   Labs zaczynają pracę nad systemem operacyjnym
%   \colorhref{https://en.wikipedia.org/wiki/Multics}{\textsc{multics}}
%   (ang.~\textit{MULTiplexed Information and Computing Service}). Prace
%   nas systemem \textsc{multics} okazały~się wyjątkowo skomplikowane i~sam
%   projekt odniósł jedynie częściowy sukces. Z~uwagi na trudności
%   przy tworzeniu tego systemu Bell Labs wycofały~się z~projektu
%   w~roku $1969$.

% \end{frame}
% % ######################################





% % ######################################
% \begin{frame}
%   \frametitle{Ken Thompson i~jego współpracownicy}


%   Jednym z~ludzi, który pracował z~ramienia Bell Labs nad \textsc{multics}em
%   jest \colorhref{https://en.wikipedia.org/wiki/Ken\_Thompson}
%   {Ken Thompson}. Po wycofaniu~się Bell Labs z~tego projektu, Thompson
%   poszukując nowych problemów do rozwiązania, zaczyna pisać programy na
%   znajdujący~się w~Bell Labs i~wówczas już rzadko używany komputer
%   \colorhref{https://en.wikipedia.org/wiki/PDP-7}{PDP-$7$}.
%   W~trackie pracy nad nimi w~pewnym momencie dochodzi do~wniosku,
%   że~niewiele mu brakuje do stworzenia systemu operacyjnego. W~ten sposób
%   powstaje system operacyjny \textsc{unix}, którego pierwsza wersja
%   jest gotowa na jesień $1969$~roku.

%   Sama nazwa „\textsc{unix}”,
%   jak zgadzają~się wszyscy biorący udział w~jego tworzeniu, jest
%   nawiązaniem, i~to dość złośliwym do nazwy „\textsc{multics}”.
%   Związane jest to z~powszechną krytyką, że~\textsc{multics} próbował
%   osiągnąć zbyt wiele rzeczy na raz i~to nie wyszło mu na zdrowie,
%   a~nowy system będzie skupiał~się na robieniu kilku rzeczy dobrze. Stąd
%   „multi” zmieniono na „unit” i~powstała nazwa \textsc{unix}, por.
%   str.~.

%   Do Thompsona w~pracy nad \textsc{unix}em dołącza wielu wybitnych
%   informatyków, pracujących wtedy w~Bell Labs, jak

% \end{frame}
% % ######################################





% % ######################################
% \begin{frame}
%   \frametitle{Ken Thompson i~jego współpracownicy}


%   System \textsc{unix} jest pierwszym bezpośrednim przodkiem używanego
%   przez nas systemu GNU/Linux, jednak droga od jednego do~drugiego
%   jest dość długa i~wyboista. Zdjęcie \eqref{} przedstawia Kena Thompsona
%   i~Denisa Ritchiego pracujących na komputerze
%   \colorhref{https://en.wikipedia.org/wiki/PDP-11}{PDP-$11$}, z~użyciem
%   dalekopisu Teletype-33???. Współtwórca \textsc{unix}a Brian W.~Kernighan
%   opatrzył to zdjęcie takim komentarzem.

% \end{frame}
% % ######################################





% % ######################################
% \begin{frame}
%   \frametitle{Dlaczego?}


%   Jak Państwo widzą,

% \end{frame}
% % ######################################





% % ######################################
% \begin{frame}
%   \frametitle{To tylko część tej historii}


%   To co przedstawiliśmy tutaj jest tylko drobnym wycinkiem historii, która
%   doprowadziła do powstania używanych przez nas systemów operacyjnych,
%   takich jak GNU/Linux.

%   Jeśli ktoś chce lepiej poznać tę historię to powinien w~pierwszym rzędzie
%   sięgnąć do krótkiej i~bardzo pięknej książki Briana W.~Kernighan
%   \textit{Jak Unix tworzył historię}
%   \parencite{Kernighan-Jak-Unix-tworzyl-historie-Pub-2021}, która była
%   już przez nas wielokrotnie cytowana. Jak sama jej tytuł wskazuje, książka
%   ta dotyczy głównie systemu \textsc{unix}, jego powstania i~rozwoju. Jednak
%   w~ostatniej części, głównie w~rozdziałach ??? i~??? opisuje drogę, jaka
%   wiodła od~\textsc{unix}, do~systemu GNU/Linux. Cała zaś książka jest
%   kopalnią informacji na temat, tego kiedy, gdzie i~dlaczego powstały znane
%   nam dziś części naszych systemów operacyjnych.

% \end{frame}
% % ######################################



% % ##################
% \jagiellonianendslide{?????}
% % ##################










% ####################################################################
% ####################################################################
% Bibliography

\printbibliography





% ####################################################################
% End of the document

\end{document}
